%%%%%%%%%%%%%%%%%%%%%%%%%%%%%%%%%%%%%%%%%%%%%%%%%%%%%%%%%%%%%%%%%%%%%%%%%%%%%%%%
\section{Predictive Retention and Performance Trade-offs (RQ3)}
\label{sec:results:rq3}
% Answers \ref{sec:intro:rq3}. The domain-informed comparison is \ref{sec:results:rq4}.

\subsection{Accuracy Degradation}
\label{sec:results:rq3:degradation}
How much accuracy survives training and testing under the same reduction, against the RQ1
full-graph baseline. RMSE, $R^2$ and Spearman correlation are reported together throughout,
because a reduction can hold mean error steady while destroying explained variance
(\ref{sec:method:experiment:metrics}).

\begin{figure*}[htbp]
    \centering
    \includegraphics[width=\textwidth]{\resultsfigs/rq3_accuracy}
    \caption[Matched-state accuracy]{\textbf{PARTLY FABRICATED.} RMSE, $R^2$ and Spearman
    correlation under the reduction each model was trained with, against the full-graph
    baseline. The three metrics disagree about the ranking. The summarization rows are
    invented, since no summarization configuration has been trained.}
    \label{fig:rq3_accuracy}
\end{figure*}

\begin{table*}[htbp]
\centering
\caption{RQ3 matched-state accuracy: models trained and tested under the same reduction. RMSE, $R^2$ and Spearman are reported together throughout -- a reduction can hold mean error steady while destroying explained variance.}
\label{tab:rq3_retention}
\small
\resizebox{\textwidth}{!}{%
\begin{tabular}{lllrrrrr}
\toprule
\textbf{Method} & \textbf{Family} & \textbf{Kind} & \textbf{Smooth L1} & \textbf{RMSE} & \textbf{$R^2$} & \textbf{Spearman} & \textbf{$R^2$ retained} \\
\midrule
Full graph & Baseline & Generic & 0.014 & 0.043 & 0.343 & 0.185 & 1.000 \\
Random hashing & Partition & Generic & 0.038 & 0.075 & -0.981 & -0.183 & -2.861 \\
METIS & Partition & Generic & 0.055 & 0.107 & -3.062 & 0.237 & -8.932 \\
Level slicing & Partition & Domain & 0.018 & 0.057 & -0.142 & 0.119 & -0.414 \\
Span-weighted METIS & Partition & Domain & 0.034 & 0.053 & 0.008 & 0.263 & 0.024 \\
Random edge dropout & Sparsification & Generic & 0.016 & 0.047 & 0.215 & 0.194 & 0.627 \\
Spanning forest & Sparsification & Generic & 0.016 & 0.049 & 0.138 & 0.282 & 0.404 \\
PageRank pruning & Sparsification & Generic & 0.019 & 0.053 & -0.018 & 0.367 & -0.051 \\
AND-gate only & Sparsification & Domain & 0.018 & 0.055 & -0.064 & -0.025 & -0.188 \\
{[TODO/FAKE] Random within-type} & Summarization & Generic & -- & 0.060 & 0.100 & 0.100 & 0.292 \\
{[TODO/FAKE] ConvMatch} & Summarization & Generic & -- & 0.050 & 0.200 & 0.200 & 0.583 \\
{[TODO/FAKE] Cone coarsening} & Summarization & Domain & -- & 0.045 & 0.300 & 0.300 & 0.875 \\
{[TODO/FAKE] MFFC coarsening} & Summarization & Domain & -- & 0.040 & 0.350 & 0.350 & 1.021 \\
{[TODO/FAKE] Colour refinement (exact)} & Summarization & Domain & -- & 0.043 & 0.343 & 0.185 & 1.000 \\
\bottomrule
\end{tabular}
}
\\[4pt]\begin{minipage}{\linewidth}\footnotesize\emph{Rows marked [TODO/FAKE] are invented. No summarization configuration has been trained or measured. Needs: sbatch src/shell/precompute\_summarization.sh for the offline pass, then one training run and one src/test.py pass per method. Random within-type merging is not implemented at all.}\end{minipage}
\end{table*}


No reduction beats the full graph on explained variance. Random edge dropout retains the
most, $R^2 = 0.215$ against the baseline's $0.343$, and the two METIS variants fall to
$-3.062$ and $0.008$. Ranking behaves differently: PageRank pruning reaches a Spearman
correlation of $0.367$ against the baseline's $0.185$, so five of the eight measured
reductions order circuits \emph{better} than the unreduced model does while explaining less
of their variance. That dissociation is the RQ3 result, and it is the reason the two metrics
are never reported apart.

\begin{figure}[htbp]
    \centering
    \includegraphics[width=\columnwidth]{\resultsfigs/rq3_metric_divergence}
    \caption[RMSE against explained variance]{RMSE against $R^2$, one point per
    configuration. Several configurations hold RMSE close to the baseline while $R^2$ falls
    below zero, which is the same dissociation \citet{cts_bench2026} report on a coarsened
    EDA graph.}
    \label{fig:rq3_metric_divergence}
\end{figure}

Holding error steady while losing all explanatory power is what a model does when it
collapses onto the label's marginal distribution, and on a corpus that is 49\% exactly zero
that collapse is cheap.

\begin{figure*}[htbp]
    \centering
    \includegraphics[width=\textwidth]{\resultsfigs/rq3_by_tier}
    \caption[Matched-state accuracy by tier]{Every configuration split across the two stored
    tiers. The full-graph baseline scores $R^2 = -0.038$ on tier-0 base graphs against
    $+0.358$ on tier-1: its absolute error on tier~0 is smaller (RMSE 0.025 against 0.047),
    but tier-0 labels have so much less spread that the model fails to beat their mean.}
    \label{fig:rq3_by_tier}
\end{figure*}

PageRank pruning is the exception, scoring $R^2 = 0.204$ and $\rho = 0.521$ on tier~0,
better than the unreduced baseline on that tier. A reduction outperforming the full graph on
a subgroup is the first sign that the reduction is acting as a regulariser rather than as a
loss of information.

\begin{figure*}[htbp]
    \centering
    \includegraphics[width=\textwidth]{\resultsfigs/rq3_by_source}
    \caption[Matched-state accuracy by source script]{The same partition taken by the
    synthesis script that produced each graph, which separates the corpus far more sharply
    than the tier does (\ref{fig:rq1_variance_decomposition}). Every configuration degrades
    from Syn4-derived graphs to C2RS-derived ones.}
    \label{fig:rq3_by_source}
\end{figure*}

\begin{table*}[htbp]
\centering
\caption{Matched-state accuracy partitioned two ways: by dataset tier, and by the synthesis script that produced the graph. Both partition the same test split, so the columns within each cut are disjoint. The source script separates the corpus far more sharply than the tier does.}
\label{tab:rq3_by_subgroup}
\small
\resizebox{\textwidth}{!}{%
\begin{tabular}{lrrrrrrrrrrrr}
\toprule
\textbf{Method} & \textbf{$R^2$ C2RS} & \textbf{$R^2$ Deepsyn} & \textbf{$R^2$ Syn4} & \textbf{$R^2$ tier0} & \textbf{$R^2$ tier0 base} & \textbf{$R^2$ tier1} & \textbf{$\rho$ C2RS} & \textbf{$\rho$ Deepsyn} & \textbf{$\rho$ Syn4} & \textbf{$\rho$ tier0} & \textbf{$\rho$ tier0 base} & \textbf{$\rho$ tier1} \\
\midrule
Full graph & -1.039 & -0.097 & 0.143 & -0.038 & -0.038 & 0.358 & -0.153 & 0.167 & 0.475 & 0.044 & 0.044 & 0.234 \\
Random hashing & -127.354 & -28.352 & -0.766 & -4.012 & -4.012 & -0.823 & -0.433 & -0.319 & 0.024 & -0.298 & -0.298 & -0.156 \\
METIS & -736.099 & -164.691 & 0.954 & -20.880 & -20.880 & -2.046 & 0.030 & 0.208 & 0.758 & 0.225 & 0.225 & 0.241 \\
Level slicing & -0.111 & -0.068 & -0.558 & -0.172 & -0.172 & -0.154 & 0.102 & 0.193 & 0.140 & 0.021 & 0.021 & 0.149 \\
Span-weighted METIS & -55.986 & -14.553 & 0.074 & -1.116 & -1.116 & 0.064 & -0.235 & 0.030 & 0.643 & -0.050 & -0.050 & 0.370 \\
Random edge dropout & -0.649 & 0.071 & -0.045 & -0.065 & -0.065 & 0.223 & 0.125 & 0.118 & 0.610 & 0.026 & 0.026 & 0.239 \\
Spanning forest & -0.342 & -0.006 & -0.153 & -0.108 & -0.108 & 0.143 & -0.133 & 0.289 & 0.786 & 0.021 & 0.021 & 0.368 \\
PageRank pruning & -3.507 & -0.626 & -0.390 & 0.204 & 0.204 & -0.043 & 0.418 & 0.194 & 0.496 & 0.521 & 0.521 & 0.312 \\
AND-gate only & -0.178 & -0.070 & -0.444 & -0.169 & -0.169 & -0.071 & -0.427 & -0.154 & 0.506 & -0.504 & -0.504 & 0.129 \\
\bottomrule
\end{tabular}
}
\end{table*}


The partitioners degrade catastrophically on the C2RS-derived subgroup, METIS reaching
$R^2 = -736$, where the labels sit near zero and it predicts high optimizability regardless.
A pooled score cannot show that, and the source script is the axis on which it hides.
\clearpage

\subsection{Ranking Preservation}
\label{sec:results:rq3:ranking}
Spearman correlation under reduction. Ranking can survive when absolute error degrades,
which is the property the script-selection use case actually needs.

\begin{figure*}[htbp]
    \centering
    \includegraphics[width=\textwidth]{\resultsfigs/rq3_stratified_matched}
    \caption[Matched-state ranking under label strata]{The matched-state result re-scored on
    four subsets of the same test split, from the persisted per-graph predictions and with
    no retraining. The strata answer different questions rather than being more or less
    trustworthy: excluding the two fixed-point designs selects on a design property,
    restricting to $y > 0$ selects on the outcome variable so its $R^2$ is not comparable to
    the pooled column, and the Syn4 column is the easiest stratum.}
    \label{fig:rq3_stratified_matched}
\end{figure*}

\begin{table*}[htbp]
\centering
\caption{Matched-state accuracy re-scored on four subsets of the same test split, from the persisted per-graph predictions. The pooled label is 49\% exactly zero, so a pooled $R^2$ is carried by a minority of graphs; a conclusion that survives all four columns is a property of the reduction rather than of the label distribution.}
\label{tab:rq3_stratified}
\small
\resizebox{\textwidth}{!}{%
\begin{tabular}{lrrrrrrrr}
\toprule
\textbf{Method} & \textbf{$R^2$ All test graphs} & \textbf{$R^2$ Designs Orchestrate can touch} & \textbf{$R^2$ Graphs with $y > 0$} & \textbf{$R^2$ Syn4-derived graphs only} & \textbf{$\rho$ All test graphs} & \textbf{$\rho$ Designs Orchestrate can touch} & \textbf{$\rho$ Graphs with $y > 0$} & \textbf{$\rho$ Syn4-derived graphs only} \\
\midrule
Random edge dropout & 0.215 & 0.192 & 0.097 & -0.045 & 0.194 & 0.225 & 0.427 & 0.610 \\
Spanning forest & 0.138 & 0.113 & 0.007 & -0.153 & 0.282 & 0.300 & 0.524 & 0.786 \\
Span-weighted METIS & 0.008 & -0.021 & 0.125 & 0.074 & 0.263 & 0.040 & 0.364 & 0.643 \\
PageRank pruning & -0.018 & -0.047 & -0.155 & -0.390 & 0.367 & 0.179 & 0.305 & 0.496 \\
AND-gate only & -0.064 & -0.095 & -0.230 & -0.444 & -0.025 & 0.179 & 0.312 & 0.506 \\
Level slicing & -0.142 & -0.175 & -0.320 & -0.558 & 0.119 & 0.548 & 0.484 & 0.140 \\
Random hashing & -0.981 & -0.061 & -0.276 & -0.766 & -0.183 & 0.134 & 0.433 & 0.024 \\
METIS & -3.062 & -3.181 & -1.997 & 0.954 & 0.237 & 0.086 & 0.276 & 0.758 \\
\bottomrule
\end{tabular}
}
\\[4pt]\begin{minipage}{\linewidth}\footnotesize\emph{Strata: all test graphs; excluding the two designs on which Orchestrate is at a fixed point (16384, 8192, max y = 0.0003 over ~32,000 graphs); graphs with a non-zero target; and graphs derived from Syn4, the one source script that leaves substantial optimizability behind.}\end{minipage}
\end{table*}


Two conclusions hold across all four strata: no reduction beats the full graph on $R^2$, and
sparsification beats partitioning. Two do not. Random hashing's $R^2$ moves from $-0.98$ to
$-0.06$ once the fixed-point designs are excluded, and METIS is worst overall yet best on
Syn4-derived graphs, which is the signature of a model that collapsed onto predicting high
optimizability rather than of a failure of partitioning.

\begin{figure*}[htbp]
    \centering
    \includegraphics[width=\textwidth]{\resultsfigs/rq3_stratified_cross}
    \caption[Cross-state ranking under label strata]{The same stratified treatment applied to
    the full-graph evaluation pass, in which each model is queried on unreduced graphs.
    Ranking separates sharply from error: level slicing reaches a Spearman correlation of
    0.601 on full graphs, above the full-graph baseline's 0.185, while AND-gate-only reaches
    $-0.467$ and orders circuits backwards.}
    \label{fig:rq3_stratified_cross}
\end{figure*}

Neither of those two is visible in $R^2$ alone, and both are large enough to change which
method a practitioner would pick. RQ5 takes the cross-state result up as a question in its
own right (\ref{sec:results:rq5:ranking}).
\clearpage

\subsection{Accuracy Loss Attributable to Reduction}
\label{sec:results:rq3:attribution}
What each method costs and what structural damage explains it. The comparison is made at
matched compression, since otherwise ``which method is best'' confounds a genuinely better
method with one that simply reduced less.

\begin{figure*}[htbp]
    \centering
    \includegraphics[width=\textwidth]{\resultsfigs/rq3_retention_vs_compression}
    \caption[Accuracy against achieved compression]{Matched-state $R^2$ against edge and
    node retention, one point per configuration. Without this plane a ranking of methods
    confounds the quality of a heuristic with the amount it happened to remove.}
    \label{fig:rq3_retention_vs_compression}
\end{figure*}

The plane is what makes the confound concrete: AND-gate-only keeps 73.3\% of edges against
spanning forest's 58.1\%, so their raw accuracies are not comparable and the formal paired
version of the comparison belongs to RQ4 (\ref{sec:results:rq4}).

\begin{figure*}[htbp]
    \centering
    \includegraphics[width=\textwidth]{\resultsfigs/rq3_accuracy_retained}
    \caption[Accuracy retained under reduction]{Each metric as a fraction of the full-graph
    baseline, one row per method, so that the figure lines up directly against the efficiency
    ranking of~\ref{sec:results:rq2:ranking}.}
    \label{fig:rq3_accuracy_retained}
\end{figure*}

Expressing accuracy as a retained fraction rather than an absolute puts it on the same
footing as the memory and time savings of RQ2, which is what the Pareto analysis below
needs.

\begin{figure*}[htbp]
    \centering
    \includegraphics[width=\textwidth]{\resultsfigs/rq3_by_target_bin}
    \caption[All configurations by optimizability band]{Every configuration as a function of
    how optimizable the circuit truly is. Level slicing predicts $0.0000$ in every band, a
    fully collapsed model that the pooled metrics do not expose. METIS over-predicts
    everywhere, which is why its pooled $R^2$ is $-3.06$, yet that same miscalibration makes
    it the \emph{most} accurate configuration on the highly optimizable circuits: MAE
    $0.046$ above $y = 0.20$ against the unreduced baseline's $0.180$.}
    \label{fig:rq3_by_target_bin}
\end{figure*}

\begin{table*}[htbp]
\centering
\caption{Mean prediction and mean absolute error within bands of the true optimizability. $R^2$ is not reported: inside a narrow band there is almost no label variance to divide by. Compare each method's predicted mean against the true mean in the first row. A row that stays flat across the bands is a model that has collapsed to a constant.}
\label{tab:rq3_by_target_bin}
\small
\resizebox{\textwidth}{!}{%
\begin{tabular}{lrrrrrrrrrrrr}
\toprule
\textbf{Method} & \textbf{pred $y = 0$} & \textbf{pred (0, .01]} & \textbf{pred (.01, .05]} & \textbf{pred (.05, .10]} & \textbf{pred (.10, .20]} & \textbf{pred $> .20$} & \textbf{MAE $y = 0$} & \textbf{MAE (0, .01]} & \textbf{MAE (.01, .05]} & \textbf{MAE (.05, .10]} & \textbf{MAE (.10, .20]} & \textbf{MAE $> .20$} \\
\midrule
(true mean in band) & 0.000 & 0.002 & 0.021 & 0.065 & 0.161 & 0.226 & -- & -- & -- & -- & -- & -- \\
Full graph & 0.004 & 0.003 & 0.005 & 0.006 & 0.052 & 0.046 & 0.004 & 0.003 & 0.016 & 0.059 & 0.109 & 0.180 \\
Random hashing & 0.042 & 0.007 & 0.002 & 0.003 & 0.008 & 0.012 & 0.042 & 0.007 & 0.019 & 0.062 & 0.153 & 0.214 \\
METIS & 0.047 & 0.075 & 0.100 & 0.109 & 0.171 & 0.180 & 0.047 & 0.073 & 0.093 & 0.083 & 0.020 & 0.046 \\
Level slicing & 0.000 & 0.000 & 0.000 & 0.000 & 0.000 & 0.000 & 0.000 & 0.002 & 0.021 & 0.065 & 0.161 & 0.226 \\
Span-weighted METIS & 0.030 & 0.032 & 0.036 & 0.016 & 0.056 & 0.052 & 0.030 & 0.029 & 0.027 & 0.049 & 0.105 & 0.174 \\
Random edge dropout & 0.004 & 0.005 & 0.006 & 0.006 & 0.023 & 0.047 & 0.004 & 0.004 & 0.015 & 0.059 & 0.138 & 0.179 \\
Spanning forest & 0.002 & 0.002 & 0.004 & 0.003 & 0.024 & 0.029 & 0.002 & 0.003 & 0.017 & 0.062 & 0.137 & 0.197 \\
PageRank pruning & 0.007 & 0.008 & 0.008 & 0.020 & 0.010 & 0.013 & 0.007 & 0.007 & 0.013 & 0.045 & 0.151 & 0.213 \\
AND-gate only & 0.002 & 0.001 & 0.000 & 0.000 & 0.006 & 0.009 & 0.002 & 0.002 & 0.021 & 0.065 & 0.155 & 0.217 \\
\bottomrule
\end{tabular}
}
\end{table*}


If the use case is finding circuits worth optimizing rather than scoring all of them, the
pooled ranking inverts and the worst method on $R^2$ becomes the best one. The banded view
is the only place in the chapter where a fully collapsed model is distinguishable from a
merely weak one.

\begin{figure}[htbp]
    \centering
    \includegraphics[width=\columnwidth]{\resultsfigs/h1_receptive_field}
    \caption[Effective receptive field]{\textbf{ENTIRELY FABRICATED. NO SUMMARIZATION DATA
    EXISTS.} Mean 4-hop fanin-cone size after reduction relative to before. The metric is
    specified in~\ref{sec:method:experiment:metrics:reduction} and is not implemented.}
    \label{fig:h1_receptive_field}
\end{figure}

\begin{figure}[htbp]
    \centering
    \includegraphics[width=\columnwidth]{\resultsfigs/h1_path_contraction}
    \caption[H1: receptive field against accuracy]{\textbf{ENTIRELY FABRICATED. NO
    SUMMARIZATION DATA EXISTS.} The form the H1 test would take: if coarsening helps by
    contracting paths, accuracy should rise with the receptive-field ratio. Both axes are
    invented.}
    \label{fig:h1_path_contraction}
\end{figure}

Hypothesis~H1 predicts that summarization retains more than sparsification at matched
compression, because contracting a path shortens the distance a signal must travel while
removing an edge can only impede it. Testing it needs both the receptive-field metric and at
least one trained summarization configuration, and neither exists, so H1 is asserted rather
than tested and must not be reported as evidenced.
\clearpage

\subsection{The Pareto Front (Trade-Off Analysis)}
\label{sec:results:rq3:pareto}
Accuracy against the memory and speed gains of RQ2, on one pair of axes, which is where RQ2
and RQ3 are answered jointly rather than separately.

\begin{figure*}[htbp]
    \centering
    \includegraphics[width=\textwidth]{\resultsfigs/rq3_pareto}
    \caption[Pareto front]{Accuracy against efficiency, with the dominance frontier drawn.
    The full-graph baseline sits on the same axes at zero saving, so ``no reduction''
    competes directly against every reduction rather than sitting outside the comparison.}
    \label{fig:rq3_pareto}
\end{figure*}

\begin{table}[htbp]
\centering
\caption{Accuracy against efficiency, ranked by matched-state $R^2$.}
\label{tab:rq3_pareto}
\small
\begin{tabular}{lrrrr}
\toprule
\textbf{Method} & \textbf{$R^2$} & \textbf{Spearman} & \textbf{VRAM saving (\%)} & \textbf{Time saving (\%)} \\
\midrule
Full graph & \textbf{0.343} & 0.185 & 0.000 & 0.000 \\
Random edge dropout & \underline{0.215} & 0.194 & 1.489 & 0.229 \\
Spanning forest & 0.138 & \underline{0.282} & 2.142 & -0.101 \\
Span-weighted METIS & 0.008 & 0.263 & 0.357 & 22.412 \\
PageRank pruning & -0.018 & \textbf{0.367} & \textbf{14.492} & 19.145 \\
AND-gate only & -0.064 & -0.025 & \underline{11.257} & \textbf{28.893} \\
Level slicing & -0.142 & 0.119 & 1.487 & 18.291 \\
Random hashing & -0.981 & -0.183 & 3.100 & 20.996 \\
METIS & -3.062 & 0.237 & 0.279 & \underline{24.284} \\
\bottomrule
\end{tabular}
\\[4pt]\begin{minipage}{\linewidth}\footnotesize\emph{\textbf{Bold} marks the best value in a column and \underline{underline} the second best, over measured rows only. Lower is better for error, memory, time and offline cost; higher is better for $R^2$, correlation and savings.}\end{minipage}
\end{table}


The baseline being on the frontier is the outcome to read first: on this corpus, at this
scale, no reduction is worth its accuracy cost if explained variance is the objective. The
argument for reduction has to be made either on a subgroup, on ranking, or at a scale where
the full graph does not fit at all.

\begin{figure}[htbp]
    \centering
    \includegraphics[width=\columnwidth]{\resultsfigs/rq3_val_test_gap}
    \caption[Validation-to-test gap]{The same configurations scored on validation and on
    test. Under a design-level split these are different circuits rather than two samples of
    one population, so the size of this gap bounds how much of the validation-selected model
    choice transfers.}
    \label{fig:rq3_val_test_gap}
\end{figure}

Every configuration on the frontier above was selected on validation, so the gap is a direct
bound on how much of that selection survives the move to held-out designs.
\clearpage
